
%(BEGIN_QUESTION)
% Copyright 2015, Tony R. Kuphaldt, released under the Creative Commons Attribution License (v 1.0)
% This means you may do almost anything with this work of mine, so long as you give me proper credit

Les utvalgte deler av instruksjonsmanualen for vernereléet ``SEL-387A Current Differential Relay'' (dokument SEL-387A Instruction Manual, januar 2014), og svar på følgende spørsmål:

\vskip 10pt

Nevn vanlige bruksområder for denne relémodellen.

\vskip 10pt

Som mange andre moderne digitale vernereléer kan modell 387A levere mer enn én ANSI/IEEE-vernefunksjon. Oppgi noen funksjoner dette reléet tilbyr utover 87 (differensialvern).

\vskip 10pt

Figur 2.7 på side 2-10 viser en eksempelapplikasjon der 387A beskytter en transformator. Identifiser viklingskonfigurasjonen i krafttransformatoren (f.eks. Delta/Stjerne), samt konfigurasjonen av primær- og sekundærstrømtransformatorene i tilkoblingen mot reléet. Hvordan skiller CT-koblingen seg fra tradisjonelle (elektromekaniske) differensialreléer?

\vskip 10pt

Basert på krafttransformator-konfigurasjonen i figur 2.7: hvor stor faseforskyvning finnes mellom primær- og sekundærfaser?

\vskip 10pt

En sentral problemstilling ved implementering av differensialvern på krafttransformatorer er {\it faseforskyvningskompensasjon} fra primær til sekundær når transformatoren er Stjerne-Delta eller Delta-Stjerne. I elektromekaniske reléer løses dette ved ulik CT-kobling på primær- og sekundærside (dvs. Delta-koblede CT-er på transformatorens stjerneside, og Stjerne-koblede CT-er på transformatorens deltaside). Digitale 87-reléer som modell 387A tilbyr derimot ``connection compensation'', basert på digitale beregninger i stedet for fysisk kobling. Side 3-17 til 3-21 beskriver denne funksjonen i modell 387A. Les denne delen og forklar med egne ord hvordan connection compensation virker for de to eksemplene i figur 3.9 og 3.10.

\vskip 20pt \vbox{\hrule \hbox{\strut \vrule{} {\bf Forslag til sokratisk diskusjon} \vrule} \hrule}

\begin{itemize}
\item{} Viklingskompensasjonsparameteren må settes til en bestemt verdi på stjernesiden av krafttransformatoren hvis CT-ene på samme side også er stjernekoblet. Forklar hvorfor, og hva {\tt WnCTC}-verdien må settes til.
\item{} Studer fasordiagrammene for {\tt WnCTC}-eksemplene på side 3-19 (figur 3.9) og 3-20 (figur 3.10). Hvilke fasordiagrammer viser faseforskyvningen i krafttransformatoren? Hvordan representeres CT-faseforskyvning? Hva representerer den stiplede linjen? Hvordan passer {\tt WnCTC}-verdiene inn i diagrammene?
\end{itemize}

\underbar{file i00824}
%(END_QUESTION)





%(BEGIN_ANSWER)

\noindent
{\bf Delsvar:}

\vskip 10pt

I figur 2.7 (side 2-10) ser vi at 387A-reléet beskytter en spenningsnedtrappende {\it autotransformator} med stjernevinding. Både primær- og sekundær-CT-er er stjernekoblet, som er vanlig for digitale 87-reléer.

%(END_ANSWER)





%(BEGIN_NOTES)

Modell 387A kan brukes til vern av transformatorer, roterende maskiner, reaktorer og andre flerterminale kraftsystemkomponenter (side 1-1). Noen applikasjoner er vist på side 1-4.

\vskip 10pt

I tillegg til differensialvern (87) tilbyr dette reléet momentant og tidsforsinket overstrømsvern (50/51) samt retningsbestemt overstrømsvern (67). Figur 1.1 på side 1-2 viser dette tydelig.

\vskip 10pt

I figur 2.7 (side 2-10) ser vi at 387A-reléet beskytter en spenningsnedtrappende {\it autotransformator} med stjernevinding. Både primær- og sekundær-CT-er er stjernekoblet, slik som ofte brukes med digitale 87-reléer. Selv om krafttransformatoren hadde Stjerne-Delta eller Delta-Stjerne, kunne CT-ene på begge sider fortsatt vært stjernekoblet fordi faseforskyvningskompensasjonen gjøres digitalt i reléet. Elektromekaniske 87-reléer uten matematisk kompensasjon krevde komplementær CT-kobling (dvs. Delta-CT på stjernesiden av krafttransformatoren, og omvendt).

\vskip 10pt

Autotransformatoren med Stjerne-Stjerne-konfigurasjon mellom primær og sekundær gir ingen faseforskyvning mellom inn- og utgang.

\vskip 10pt

Kompensasjon for viklingsforbindelse er en funksjon som aktiveres/deaktiveres med innstillingen {\tt ICOM} (side 3-9). Når funksjonen er aktiv, brukes parameterne {\tt W1CTC} og {\tt W2CTC} (CT-kompensasjon for vikling 1 og 2), som tar heltall fra 1 til 12 og angir ekstra faseforskyvning som legges til CT-signalene for den aktuelle viklingen. Verdiene følger ``klokketimer'', der hver time tilsvarer 30 graders faseforskyvning. Dermed betyr ``3'' +90 grader, ``6'' +180 grader og ``10'' +300 grader. Hvis verdien ``12'' settes, fjernes nullsekvensstrømmer fra den CT-gruppen, noe som er nødvendig dersom CT-ene er stjernekoblet på stjernesiden av krafttransformatoren.

%INDEX% Electric power systems: protective relays (differential)
%INDEX% Reading assignment: SEL model 387A current differential protective relay

%(END_NOTES)
